\section{Graph consequences of path and segment costs}\label{met:d-section}
The next four models change what a tree representation costs: a sum of
Hilbert and path costs, their Euclidean combination, a charge at each
segment's start, or a restriction to incomparable segments. These are
different norms. Within one norm, several estimates can also give
different bounds. We first extract their squared losses, then retain
the actual-child or qualitative conclusions when the method supplies
different information. All intrinsic inputs come from the companion on
bounded tree potentials and path costs; the graph deductions use the
transfers proved above.

\subsection{Linear path cost}\label{met:linear-section}
Recall the linear path-cost space $X_{\rm lin}$ from
Section~\ref{dir:layers}. Its component norm is the infimum of
$\|h\|_2+\|\mu\|_1$ over decompositions into a Hilbert vector and
root paths, including the component root. For later comparison write
$J_n\mu=\sum_s\mu(s)b_s$ for the root-path sum on
$T_n=\mathbb N^{\le n}$.
The finite ancestor-closed coordinate projections are contractive,
finite supports are dense, and this space and its dual are separable and
reflexive. The squared estimates below and the weakly null estimate in
Section~\ref{met:weak-linear} give lower bounds for this same norm;
the latter will give the strongest coefficient displayed here.

The positive-flow argument gives, for $Pp=p$, $Py=0$
\cite[Theorem~5.7]{TreePotentials},
\begin{equation}\label{met:linear-flow-input}
 \|y\|^2\le C^2(\|p+y\|^2+\|p-y\|^2-2\|p\|^2),
 \qquad C=6(\sqrt2+1).
\end{equation}
A separate trimming argument \cite[Theorem~5.10]{TreePotentials}
gives the same inequality with $C^2$
replaced by $242$. Since each endpoint is at most $R$, these are
\eqref{met:square-input} with $a=1/(2C^2)$ and $a=1/484$,
respectively. Proposition~\ref{met:tail-to-lens} and
Theorem~\ref{met:squared-transfer} therefore give
\begin{equation}\label{met:linear-flow-graph}
 D^2\ge1+\frac{k}{2C^2},\qquad
 D^2\ge1+\frac{k}{484},
\end{equation}
using the corresponding proof in each case. In upper-one normalization,
the first argument says $r_j^2\le r_{j+1}^2-1/(2C^2D^2)$ and hence
also retains $k\le2C^2D^2$.

The inward-clipping argument \cite[Theorem~5.4]{TreePotentials}
gives the different parameterized estimate
\begin{equation}\label{met:linear-clipping-input}
 \max_\pm\|p\pm y\|
 \ge\frac{\|p\|+c\|y\|/2}{\sqrt{1+4c^2}}
 \qquad(c>0, Pp=p, Py=0).
\end{equation}
Here is its direct actual-child application, retaining the error budget.
For a parent $x$ in the noncontracting normalization set
\[
 c=\frac1{8D},\quad Q=\sqrt{1+4c^2},\quad
 \gamma=\frac c2-(Q-1)D,\quad
 \eta=\frac{\gamma}{2(1+c/2+2Q)}.
\]
Since $Q-1\le2c^2$, we have $\gamma\ge c/4>0$.
Choose a finite head $P$ with $\|x-Px\|<\eta$, and choose the pair in
Lemma~\ref{met:finite-pair}; write $w=(z_i-z_j)/2$ and $y=w-Pw$.
The greatest of the four actual child norms is at least
\begin{align*}
 \max_\pm\|x\pm w\|
 &\ge \frac{\|Px\|+c\|y\|/2}{Q}-2\eta\\
 &\ge \frac{\|x\|-\eta+c(1-\eta)/2}{Q}-2\eta
 \ge\|x\|+\frac{\gamma}{2Q}.
\end{align*}
Selecting such a child through all $k$ generations gives
\begin{equation}\label{met:linear-clipping-graph}
 1+\frac{k\gamma}{2Q}\le D.
\end{equation}
The additive increment is attached to the clipping estimate, independently
of the positive-flow and trimming proofs.

\subsection{Weakly null pairs for the linear norm}\label{met:weak-linear}
A further form of inward clipping \cite[Theorem~5.5]{TreePotentials}
states that for every $x\in X_{\rm lin}$
and every bounded weakly null sequence $(y_j)$,
\begin{equation}\label{met:weak-input}
 \liminf_j\max_\pm\|x\pm y_j\|
 \ge\sqrt{\|x\|^2+\frac1{16}(\liminf_j\|y_j\|)^2}.
\end{equation}
This input retains a larger squared loss when we pair deviations before
applying it:
\begin{equation}\label{met:linear-weak-graph}
 D^2\ge1+k/16.
\end{equation}

For a parent with child bound $R$ in the noncontracting normalization,
the bounded sequence of deviations $z_i$ has a weakly Cauchy
subsequence. Indeed, diagonalize its evaluations on a countable
norm-dense subset of $X_{\rm lin}^*$, then use boundedness to
approximate every other functional. Relabel this subsequence and set
$w_j=(z_{2j}-z_{2j-1})/2$. It is weakly null, and sibling separation
gives $\|w_j\|\ge1$. Crossed convex averages give
\[
 x+w_j=\tfrac12((x+z_{2j})+(x-z_{2j-1})),\qquad
 x-w_j=\tfrac12((x-z_{2j})+(x+z_{2j-1})),
\]
so both norms are at most $R$. Equation~\eqref{met:weak-input} yields
$\|x\|^2+1/16\le R^2$. Taking the stage suprema and summing proves
\eqref{met:linear-weak-graph}. Alternatively, reflexivity supplies a
weak limit $z$ and one can use $(z_i-z)/2$. Separation then gives only
$\liminf\|(z_i-z)/2\|\ge1/2$, yielding the weaker
$D^2\ge1+k/64$; pairing two deviations avoids that loss.

The qualitative bottom-up mass-removal argument
\cite[Theorem~5.9]{TreePotentials} can also be transferred
without replacing it by~\eqref{met:weak-input}. It states that pure-tail
perturbations tend uniformly to zero as finite-head centers approach
the unit sphere, allowing the heads to vary. Thus it satisfies the
hypothesis of Lemma~\ref{met:bridge:sequential-lens}, which supplies
a uniform separated-lens gap. Proposition~\ref{met:backward} gives
$D^{-1}\le\rho_D^k$ for some $\rho_D<1$. This route uses the
qualitative input alone and requires no weak-limit extraction or
numerical modulus.

\subsection{Quadratic path cost}\label{met:quadratic-section}
On the same trees define
\[
 Q_n(z)=\inf_{z=h+J_n\mu}(\|h\|_2^2+\|\mu\|_1^2)^{1/2},\qquad
 X_{\rm quad}=\left(\bigoplus_{n\ge1}(\ell_2(T_n),Q_n)\right)_2.
\]
This changes the norm: the root singleton has quadratic-cost norm
$1/\sqrt2$, whereas its linear-cost norm is one. The two cancellation
estimates \cite[Theorem~7.6 and Proposition~7.10]{TreePotentials}
have the same arbitrary-displacement form, with constants $8$ and $7$:
\begin{equation}\label{met:quadratic-input}
 \|(I-P)y\|\le K\sqrt{R^2-\|p\|^2},
 \qquad Pp=p,\quad \|p\pm y\|\le R,\qquad K\in\{8,7\}.
\end{equation}
Finite ancestral heads have bounded finite-rank projections, and finite
coordinate vectors are dense. These facts permit the same packing
argument for either constant.
For an infinite $\varepsilon$-separated $R$-lens at arbitrary $x$, choose
a finite-support $p$ with $\|p-x\|<\eta$. Apply
\eqref{met:quadratic-input} to every deviation at radius $R+\eta$,
using the same finite-rank head $P$. Two head projections are arbitrarily
close. Hence
\[
 \varepsilon\le2K\sqrt{(R+\eta)^2-\|p\|^2}.
\]
Letting $\eta$ decrease to zero gives
$\|x\|^2\le R^2-\varepsilon^2/(4K^2)$. Thus
Theorem~\ref{met:squared-transfer} gives $D^2\ge1+k/K^2$.
In particular,
\begin{equation}\label{met:quadratic-graph}
 D^2\ge1+k/49.
\end{equation}
The coefficient-$8$ proof independently gives $D^2\ge1+k/64$.
Only boundedness and finite rank of $P$ entered the packing step.

The separate stopped-flow energy estimate
\cite[Proposition~7.8]{TreePotentials} has input
\[
 \|p\|^2+\|y\|^2/144
 \le\tfrac12(\|p+y\|^2+\|p-y\|^2),
 \qquad Pp=p,\quad Py=0.
\]
It gives $r_j^2\le r_{j+1}^2-1/(144D^2)$ in upper-one normalization,
and hence the quantitative necessary condition $k\le144D^2$.
The finite-head pairing, the $R+2\eta$ endpoint error, and the limiting
step are exactly those in Proposition~\ref{met:tail-to-lens}.

\paragraph{Packing at a fixed radius.}
The same arbitrary-tail input also allows a direct finite-error packing
argument. Its coefficient-$7$ version gives, for $\alpha=1/D$ and
$\delta=\alpha^2/2000$, the backward factor $1-\delta$. Indeed, if a
unit-ball lens had center norm greater than $1-\delta$, approximate it
within $\delta$ by $p$. Then $\|p\|>1-2\delta$ and the endpoint
radius is $1+\delta$. Its common tails have norm at most
\[
 7\sqrt{(1+\delta)^2-(1-2\delta)^2}
 \le7\sqrt{6\delta}<\alpha/2.
\]
Two head projections within $\alpha$ contradict separation $2\alpha$.
For the pruning estimate with coefficient $9$
\cite[Proposition~7.12]{TreePotentials}, choose instead
\begin{equation}\label{met:pruning-parameter}
 0<\gamma<1/2,\qquad
 9\sqrt{(1+\gamma)^2-(1-2\gamma)^2}<\frac{2\alpha}{3}.
\end{equation}
The same approximation bounds every tail by $2\alpha/3$; choose two
head projections within $2\alpha/3$. Their full distance is less than
$2\alpha$, a contradiction. Thus the backward factor is $1-\gamma$.
These choices retain the bounds supplied by their respective proofs.

The threshold-removal argument \cite[Proposition~7.14]{TreePotentials}
gives a different relative-radius input:
for every nonzero $x$ and $0<\delta\le1$, a single finite
head bounds all deviations satisfying
$\|x\pm y\|\le(1+\delta)\|x\|$ by
\[
 \|(I-P)y\|\le72\sqrt\delta\,\|x\|.
\]
Take $\delta=(200D)^{-2}$. If every child of a noncontracting parent
had norm at most $(1+\delta)\|x\|$, their deviations would have a
common tail bound $72\sqrt\delta D$. Two head projections can be
chosen within one, giving full distance less than
$1+144\sqrt\delta D<2$. This contradicts sibling separation. An
actual child therefore grows by the factor $1+\delta$, and
$D\ge(1+\delta)^k$.

\paragraph{A finite-codimensional input.}
The component-to-sum argument
\cite[Lemma~7.15, Proposition~7.16 and Corollary~7.17]{TreePotentials}
for the quadratic norm gives, at arbitrary unit centers, a finite-codimensional subspace
on which
\[
 \max_\pm\|x\pm y\|\le1+e
 \quad\Longrightarrow\quad \|y\|\le17(2e)^{1/6},
 \qquad0<e\le1/16.
\]
Given $D$, choose $e$ so that $17(2e)^{1/6}<1/(2D)$.
Proposition~\ref{met:quotient-transfer} then gives an actual child gain
$1+e/2$, and $D\ge(1+e/2)^k$. This application uses the proved
component-to-sum estimate, including its zero-center cases, rather than
an assertion that passing to an outer Hilbert sum automatically
preserves a midpoint gap.

\subsection{Segments with prescribed start cost}\label{met:segment-section}
For nonnegative start masses $m_s$ on $T_n$, let
$q_s=m_s+(\sum_{c\text{ child of }s}q_c^2)^{1/2}$, computed upward
from the leaves. A finite or countable signed segment representation
$z=\sum_j a_j\mathbf1_{[s_j,t_j]}$, with absolutely convergent output
at every coordinate, has start masses
$m_s=\sum_{j:s_j=s}|a_j|$ and cost $q_o$ at the root. Take the
infimum over all such representations of finite cost and complete; the outer Hilbert sum
of the resulting finite-height spaces is denoted $X_{\rm start}$.
Its intrinsic inequality is \cite[Theorem~8.4]{TreePotentials}
\begin{equation}\label{met:segment-input}
 \max_\pm\|p\pm y\|^2\ge\|p\|^2+\|y\|^2/64,
 \qquad P_Cp=p,\quad P_Cy=0,
\end{equation}
where $C$ is a finite ancestral head. Coordinate projections onto such
heads are bounded and finite supports are dense.

For a noncontracting parent $x$ let $M$ be the supremum of its child
norms. For arbitrarily small $\eta>0$, choose $p$ and the paired tail
$y$ by Lemma~\ref{met:finite-pair}, with $\|y\|\ge1-\eta$.
Then
\[
 M\ge\sqrt{\|p\|^2+(1-\eta)^2/64}-2\eta.
\]
On letting $\eta$ decrease to zero, $M^2\ge\|x\|^2+1/64$.
Apply this to every parent before taking the stage suprema. Summing
their squared losses gives
\begin{equation}\label{met:segment-graph}
 D^2\ge1+k/64.
\end{equation}
If one instead wants to follow actual children, some child has squared
norm greater than
$\|x\|^2+1/128$: otherwise the supremum would be at most this
strictly smaller number. Selecting an actual child at each generation
retains the separate bound $D^2\ge1+k/128$.
No contractivity of $I-P_C$ was used, and no child was presumed to
attain the supremum.

\subsection{Incomparable segments}\label{met:incomparable-section}
Let $Y_I$ be the completion on the disjoint union of $T_n$ for the
supremum of Euclidean sums of segment evaluations, where the admissible
segments are pairwise incomparable: every node of one segment is
incomparable with every node of another. Put $X_I=Y_I^*$. This
restriction differs from requiring the segments merely to be disjoint.
The intrinsic estimate for $X_I$ \cite[Theorem~9.5]{TreePotentials}
is
\[
 \|(I-P_E)y\|\le K\sqrt{1-\|p\|},\qquad
 P_Ep=p,\quad \|p\pm y\|\le1,\quad K=\sqrt2+\sqrt6,
\]
where $E$ is finite and ancestral and $y$ may have head coordinates.
Finite ancestral projections are contractive and finite coordinate
vectors are dense. For a pure tail, homogeneity gives
$\|y\|^2\le K^2R(R-\|p\|)$ at radius $R$.
Proposition~\ref{met:tail-to-lens} and
Theorem~\ref{met:squared-transfer} therefore give
\begin{equation}\label{met:incomparable-additive}
 k\le K^2D(D-1).
\end{equation}
At radius one the same lens estimate reads, for an infinite
$2\alpha$-separated family,
\[
 \|x\|\le1-(\alpha/K)^2.
\]
Under lower factor $\alpha=1/D$ and upper factor one, normalize each
parent lens by its child bound $R\le1$. Its separation is at least
$2\alpha/R\ge2\alpha$, so every parent bound contracts by
$\rho=1-(\alpha/K)^2$. Hence $\alpha\le\rho^k$ by
Proposition~\ref{met:backward}. This is an application to the
incomparable-segment norm itself, regardless of the equality of its
displayed tail constant with one obtained for the disjoint-segment norm.
